
\section{The global trace formula in invariant form}

The global trace formula is an equality of the form 
\begin{equation}
\sum_{\chi \in {\cal X}}J_\chi (f) =\sum_{{\cal{O}}\in {\cal C}}
J_{{\cal{O}}}(f),
\end{equation}
where $J_\chi (f)=J_\chi ^T(f)|_{T=0}$,
$J_{{\cal{O}}}(f)=J_{{\cal{O}}}^T(f)|_{T=0},$ and where
$J_\chi ^T(f)$ is given in section 9.3.3 and
$J_{{\cal{O}}}^T(f)$ is given in section 9.2.8.

We can very crudely index the terms $J_{\cal{O}}$ by the centralizer
of the semi-simple part $s$ of a $\gamma =su\in {\cal{O}}$. As a
general principle, the ``smaller'' the centralizer the easier it is
to deal (analytically) with the term. 

Let ${\cal C}(G)$ denote the subset
of elliptic regular equivalence classes in ${\cal C}$ and let
${\cal{X}}(G)$
denote the subset of $\chi \in {\cal{X}}$ of the form $\chi=(G,\rho )$ for
some cuspidal
$\rho \in \Pi _u(\overline{G({\Bbb {A}})})$.

Not all the expressions in the above trace formula (8.50)
are invariant in the sense that $J(f)=J(f^y)$ for all
$f\in {\cal{H}}(\overline{G({\Bbb {A}})}) $ and all
$y\in \overline{G({\Bbb {A}})}$. The terms $J_{{\cal{O}}}(f)$
with ${\cal{O}}\in {\cal{C}}(G)$ are invariant and the terms
$J_\chi (f)$ with $\chi \in {\cal{X}}(G)$ are invariant.
We follow \cite{A6}, \cite{A9}, \cite{A10}, in putting
the remaining terms in invariant form.

Fix a finite set of places $S$ of ${\Bbb {Q}}$ which
unless stated otherwise
will contain the archimedian place and the primes $p|2n$.
Let ${\Bbb {A}}_S=\Pi _{p\in S}{\Bbb {Q}}_p$ and embed
${\cal{H}}(\overline{G({\Bbb{A}}_S)})
\hookrightarrow {\cal{H}}(G({\Bbb {A}))}$
in the natural way. We identify $f\in {\cal{H}}(\overline{G({\Bbb{A}}_S)})$
with the $f_S\in {\cal{H}}(G({\Bbb {A}}))$ which agrees
with $f$ on $\overline{G({\Bbb{A}}_S)}$ and with the unit
element of ${\cal{H}}(\overline{G({\Bbb{Q}}_p)},K_p)$
on $\overline{G({\Bbb{Q}}_p)} $ for $p\notin S$.

Define 
\begin{equation}
I_\chi (f)=I_\chi ^S(f)=J_\chi (f),\;\;\;\;\;\;\;\;\chi \in {\cal{X}}(G),
\end{equation}
and 
\begin{equation}
I_{{\cal{O}}}(f)=I_{{\cal{O}}}^S(f)=J_{{\cal{O}}}(f),
\;\;\;\;\;\;\;\;{\cal{O}}\in {\cal C}(G),
\end{equation}
for $f\in {\cal{H}}(\overline{G({\Bbb{A}}_S)})$. One of the objects to this subsection
is to define $I_\chi (f)=I_\chi ^S(f)$,
$\chi \in {\cal{X}}-{\cal{X}}(G), $
and $I_{{\cal{O}}}(f)=I_{{\cal{O}}}^S(f)$,
${\cal{O}}\in {\cal C}-{\cal C}(G),\;$%
following \cite{A6}. Once this is done, we will use the formalism of
\cite{A6} to
deduce the global invariant trace formula in the form 
\begin{equation}
\sum_{\chi \in {\cal{X}}}I_\chi (f)=\sum_{{\cal{O}}\in {\cal C}}
I_{{\cal{O}}}(f).
\end{equation}

To do this requires, in particular, verifying for each $p\in S$ the local
assumptions of \cite{A6}, of which there are three types:

(1) \cite{A6}, Assumption 5.1, p. 30,

(2) \cite{A6}, Assumption 5.2, p. 32,

(3) the statements (7.1)-(7.6) in \cite{A6}
regarding normalized intertwining operators.
(In our case, only the statements (7.1)-(7.4) need proving since
$SL(2)/k$ is rank one.)

\subsection{A function used in conjugation}

We shall now make precise three types of statements alluded to
above which must be verified.

Let $Q=M_Q N_Q$ be a parabolic of $G$ and let
$y\in \overline{G({\Bbb{A}}_S)}$.
For $m\in \overline{M({\Bbb{A}}_S)}$, let
\begin{equation}
f_{Q,y}(m)=\delta_Q(m)^{1/2}\int_K\int_{N_Q({\Bbb{A}}_S)}
f(k^{-1}mnk)u'_Q(k,y)dndk,
\end{equation}
where
\[
u'_Q(k,y)=\int_{{\frak{a}}_Q}
\sum\Sb P \\Q\subset P\endSb
(-1)^{dim(A_Q/A_G)}
\tau_Q^P(X)\hat{\tau}_P(X+H_Q(ky))dX ,
\]
as in \cite{A6}, {3.3} (modulo a typo).
(Recall (2.7) for $\tau_Q^P,...$.)
\begin{example} Let $G=SL(2)$. Then
\[
\begin{array}{c}
u'_G(k,y)=\tau_G^G(X)\hat{\tau}_G(X+H_G(ky))=1,\\
u'_Q(k,y)=\int_{{\Bbb{R}}}
-\tau_B^B(X)\hat{\tau}_B(X+H_B(ky))
+\tau_B^G(X)\hat{\tau}_G(X+H_B(ky)) dX \\
=\int_{{\Bbb{R}}}
-\hat{\tau}_B(X+H_B(ky))+\tau_B(X)dX\\
=|H_B(ky)|.
\end{array}
\]
\end{example}

\subsection{Assumption 5.1}

First, a

\begin{definition}
Let $\phi :U\rightarrow V$ denote a continuous linear transformation between
two complete topological vector spaces. Let $U^{\prime }$ denote the space
of continuous linear functionals on $U$ (similarly for $V^{\prime }$). The 
{\bf transpose} $\phi ^{\prime }:V^{\prime }\rightarrow U^{\prime }$ is the
map defined by 
\[
\phi ^{\prime }(\lambda ):f\longmapsto \lambda (\phi (f)),\;\;\;\lambda \in
V^{\prime },\;f\in U. 
\]
\end{definition}

Let $C_N^\infty({\frak{a}}_{M,S})$ denote the topological vector space of
smooth functions on ${\frak{a}}_{M,S}$ which are supported on
\[
\{X\in {\frak{a}}_{M,S}\ |\ ||X||\leq N\},
\]
where
\begin{equation}
{\frak{a}}_{M,S}
=\{H_M(m)\ |\ m\in \overline{M({\Bbb{A}}_S)}\}.
\end{equation}

We follow \cite{A11}, section 11 (as opposed to
\cite{A6}, sections 4,5 where
the assumptions are stated), restricting ourselves to
$G=SL(2)$ and simplifying whenever possible.

 Let $N,\Gamma$ be as in section 8.3 above and let
$I_N(\overline{G({\Bbb{A}}_S)})_\Gamma$ denote the vector
space of complex valued functions $\phi$
on the tempered dual $\Pi_t(\overline{G({\Bbb{A}}_S)})$ such
that

(a) if $\pi|_K$ contains no representation in $\Gamma$ then
$\phi (\pi)=0$,

(b) if $M\subset G$ is an admissible Levi (e.g., $M=A$) and
$\sigma\in \Pi_t(\overline{M({\Bbb{A}}_S)})$ then
\begin{equation}
\phi(\sigma,X)
=\int_{{\frak{a}}_{M,S}}
\phi(\sigma,X)e^{-\Lambda (X)}d\Lambda ,\ \ \ \ X\in {\frak{a}}_{M,S},
\end{equation}
converges to a function of $X$ which belongs to
$C_N^\infty({\frak{a}}_{M,S})$.

If we fix a family of semi-norms $||...||_{M,\sigma}$
on $C_N^\infty({\frak{a}}_{M,S})$ then there are seminorms
\begin{equation}
\phi\longmapsto ||\phi(\sigma,*)||_{M,\sigma},\ \ \ \
\sigma\in \Pi_t(\overline{M({\Bbb{A}}_S)}),
\end{equation}
which make $I_N(\overline{G({\Bbb{A}}_S)})_\Gamma$
into a topological vector space.
By \cite{Ru}, ch 1, we may complete
$I_N(\overline{G({\Bbb{A}}_S)})_\Gamma$ using this family of seminorms to
obtain a Frech\'et space.

Define
\begin{equation}
I(\overline{G({\Bbb{A}}_S)})
=\lim_{\Gamma}
I(\overline{G({\Bbb{A}}_S)})_\Gamma
\end{equation}
where
\[
I(\overline{G({\Bbb{A}}_S)})_\Gamma
=\lim_N
I_N(\overline{G({\Bbb{A}}_S)})_\Gamma .
\]

In \cite{A11}, the spaces $I_{ac}(\overline{G({\Bbb{A}}_S)})$ are also
introduced, as well as the subspace of moderate functions,
which we denote by $I_{mod}(\overline{G({\Bbb{A}}_S)})$.
Actually only the non-metaplectic versions of these spaces
were introduced in \cite{A11}, their metaplectic analogs are left
to the reader.
Note that for $G=SL(2)$ or $G=A$, the diagonal maximal split torus
of $SL(2)$, we have
\[
I_{ac}(\overline{G({\Bbb{A}}_S)})=I(\overline{G({\Bbb{A}}_S)})
=I_{mod}(\overline{G({\Bbb{A}}_S)}),
\ \ \ \ 
I_{ac}(\overline{A({\Bbb{A}}_S)})
=I(\overline{A({\Bbb{A}}_S)})
=I_{mod}(\overline{A({\Bbb{A}}_S)}).
\]

Let $M\in \{A,G\}$, let ${\cal{H}}_{ac}$ be as in
\cite{A11} (see also section 8.3 above),
and let 
\[
U_M={\cal{H}}_{mod}(\overline{M({\Bbb{Q}}_p)})
={\cal{H}}(\overline{M({\Bbb{Q}}_p)}) .
\]
These are smooth, compactly supported (in our case), $K$-finite
functions. Let 
\[
V_M=I_{mod}(\overline{M({\Bbb{Q}}_p)}) =I(\overline{M({\Bbb{Q}}_p)})
\]
denote the space of functions on $\Pi _u(\overline{M({\Bbb{Q}}_p)})$ as
in \cite{A11} (and as above).
Note that these function spaces are not the same as those in \cite{A6}.
The Assumption 5.1 in \cite{A6} is needed for
the spaces ${\cal H}_{mod}$, ${\cal H}_{ac}$, ${\cal H}$.
Then the assumptions are

\begin{conjecture}
The map $\phi=\phi_M^M :U_M\rightarrow V_M$ given by
$\phi (f):\pi \longmapsto tr(\pi (f)) $ is surjective.
\end{conjecture}

The non-metaplectic analog of this
was proven in \cite{A10}, Lemma 6.1 (see also \cite{BDK}).
For a discussion of the metaplectic conjecture, see chapter 6
in the p-adic case and see
Clozel and Delorme \cite{CD} and Bouaziz \cite{Bou} in the complex case
or real case with $n=1$. We shall {\bf assume} that there is a
metaplectic analog of the Paley-Wiener theorem of
Clozel-Delorme in the real case and of
Bernstein-Deligne-Kazhdan in the $2$-adic case
with $n=2$ for the rest
of this work. (For reasons discussed in the introduction above,
this will not be needed for our main result.)

\begin{conjecture}
The image of $\phi ^{\prime }:V_M^{\prime }\rightarrow U_M^{\prime }$ is the
space of all invariant distributions on
$\overline{M({\Bbb {Q}}_p)}$.
\end{conjecture}

In \cite{A6}, Lemma 8.3, Theorem 12.1, it is shown that these imply

\begin{theorem} Assume the conjectures above.
There is a map $\phi^G_M :U_G\rightarrow V_M$ satisfying

(a) if $G=M$ then $\phi^G_M$ is as above,

(b) in general, we have
\[
\phi_M^G(f^y)=
\sum\Sb Q \\ M\subset M_Q\endSb
\phi_M^{M_Q}(f_{Q,y}),\ \ \ \ f\in U_G.
\]
In particular, $\phi_G^G(f^y)=\phi_G^G(f)$.

(c) for $f\in U_G$, $\pi\in \Pi_t(\overline{M({\Bbb{A}}_S)})$,
we have $\phi^G_M(f,\pi)={\rm tr}(I_{P_0}(\pi,f)R_M^G(\pi,P_0))$,
where $R_M^G(\pi,P_0))$ is defined in \cite{A6}, section 7
via the $(G,M)$-family
\[
R_P(\lambda, \pi,P_0))=
R_{P|P_0}(\pi)^{-1}R_{P|P_0}(\pi_\lambda),
\]
where the $R_{P|P_0}(\pi_\lambda)$ are as in section 9.4.4 below.
\end{theorem}

\begin{example} Let $G=SL(2)$. In this case,
\[
\phi^G_A(f,\pi)={\rm tr}(I_{B}(\pi,f)R_A^G(\pi,B))
\]
is a global ``weighted character'' similar to the integrand of
the expression in Theorem 307 above. The term $R_A^G(\pi,B)$
calculated in a similar fashion to $M_B^T(\pi)_\chi$
in Examples 225, 308.
\end{example}

Let $\sigma\in \Pi(\overline{A({\Bbb{A}}_S)})$ and let
\[
R_{\overline{B}|B}(\sigma)
=m(\sigma)^{-1}M_{\overline{B}|B}(\sigma)
\]
Then we may write  
\[
R_{\overline{B}|B}(\sigma_\lambda)^{-1}
R'_{\overline{B}|B}(\sigma_\lambda)|_{\lambda =0}
=-m(\sigma_\lambda)^{-1}m'(\sigma_\lambda)|_{\lambda =0}
+M_{\overline{B}|B}(\sigma_\lambda)^{-1}
M'_{\overline{B}|B}(\sigma_\lambda)|_{\lambda =0}
\]
Eventually we will want to plug this into the terms 
(9.49) (using the last equation in Example 303). 

\subsection{Assumption 5.2}

Let $M\in \{A,G\}$, $S=\{p\}$.

A complex valued function $\phi$ on
$\Pi(\overline{M({\Bbb{Q}}_p)})$ extends to the
complex Grothendieck group
$R(\Pi(\overline{M({\Bbb{Q}}_p)}))\otimes {\Bbb{C}}$ by linearity.
Thus, if $L$ is any Levi of $M$ then $\phi$ is defined on the image of
the induction map
$I_L^M: \Pi(L({\Bbb{Q}}_p))\rightarrow R(\Pi(\overline{M({\Bbb{Q}}_p)}))$.
In particular, for $\sigma\in \Pi(\overline{M({\Bbb{Q}}_p)})$ and
$\zeta\in \oplus_{v\in S}i{\frak{a}}^*_v$, $\phi(I_L^M(\sigma_\zeta))$ is
defined. Let
\begin{equation}
D_\sigma \phi(\sigma^M)=
\left\{
\begin{array}{c}
{d\over d\zeta}\phi(I_L^M(\sigma_\zeta))|_{\zeta =0},
\ \ \ \ {\rm if\ it\ exists},\\
\infty,\ \ \ \ {\rm otherwise}.
\end{array}
\right.
\end{equation}
If $p<\infty$ is a prime and
$\sigma\in \Pi(\overline{M({\Bbb{Q}}_p)})$ then define
$||\sigma||=0$. If $p=\infty$,
$\sigma\in \Pi(\overline{M({\Bbb{Q}}_p)})$,
and $\Delta$
is a Casimer operator on $\overline{M({\Bbb{Q}}_p)}$
(which is the same as the Casimer operator on
$M({\Bbb{Q}}_p)$)
then define $||\sigma||$ to be the
smallest (positive) eigenvalue of $\sigma(\Delta)$. For
$\sigma\in \Pi(\overline{M({\Bbb{Q}}_p)})$, define
\[
||\sigma||=\sup_{p\in S}||\sigma_p||.
\]

If $U_M={\cal{C}}(\overline{M({\Bbb{Q}}_p)})$, let
$V_M$ be defined as follows.
Let
\[
\begin{array}{c}
{\cal{I}}(\overline{M({\Bbb{Q}}_p)})_{\overline{K}}
=\{\phi: \Pi(\overline{M({\Bbb{Q}}_p)})\rightarrow {\Bbb{C}} \ | \\
(a) \ \phi(\pi)=0\ {\rm \ unless\ }\pi{\rm \ has\ a\ }
\overline{K\cap M({\Bbb{Q}}_p)}-{\rm \ fixed\ vector},\\
(b)\ {\rm for \ each\ Levi\ }L{\rm\ of\ }M,\\
||\phi||_{D,n}=\sup_{\sigma\in \Pi_t(\overline{L({\Bbb{Q}}_p)})}
(1+||\sigma||)^n|D_\sigma \phi(\sigma^M)|\\
{\rm is\ finite}\},
\end{array}
\]
topologized by the semi-norm topology. Let
\begin{equation}
V_M=
{\cal{I}}(\overline{M({\Bbb{Q}}_p)})
=\lim_{K}{\cal{I}}(\overline{M({\Bbb{Q}}_p)})_{\overline{K}},
\end{equation}
topologized using the direct limit topology.
See also \cite{A16}, section 3.

Then the assumptions are

\begin{conjecture}
The map $\phi =\phi_M^M :U_M\rightarrow V_M$ given by
$\phi (f):\pi \longmapsto tr(\pi (f)) $ is surjective.
\end{conjecture}

This is proven in \cite{A16}, appendix.

\begin{conjecture}
The image of $\phi ^{\prime }:V_M^{\prime }\rightarrow U_M^{\prime }$ is the
space of all tempered invariant distributions on
$\overline{M({\Bbb{Q}}_p)}$.
\end{conjecture}

When $M=G=SL(2)$ these conjectures are discussed in chapter 6, where
they are shown to be a consequence of other conjectures (!).
The $p$-adic case of these conjectures are known for
$p>2$. We shall assume them for $p=2$ and $p=\infty$.


In \cite{A6}, it was shown that these imply
\begin{theorem} Assume the two conjectures above.
There is a map
$\phi^G_M :U_G\rightarrow V_M$ satisfying

(a) if $M=G$ then $\phi_M^G$ is as above.

(b) In general, we have
\[
\phi_M^G(f^y)=
\sum\Sb Q \\ M\subset M_Q\endSb
\phi_M^{M_Q}(f_{Q,y}),\ \ \ \ f\in U_G.
\]
In particular, $\phi_G^G(f^y)=\phi_G^G(f)$.

(c) $\phi_M^G$ is given as in Theorem 308(c) above.
\end{theorem}

\subsection{Normalized intertwining operators}

Let $P,P^{\prime },P^{\prime \prime }$ be parabolics of
$G$ defined over ${\Bbb{Q}}$
with Levi $A$. Let $\chi \in \Pi _u({\overline{A({\Bbb{Q}}_p)}})$ have
representation space $V_\chi $ - a finite dimensional complex vector space
with an ${A({\Bbb {Q}}_p)}$-invariant inner product
$(...,...)_\chi $ - and let ${\cal{K}}_P(\chi )$ denote the space of all
genuine $V_\chi $-valued functions $\phi $
whose norms are square-integrable over
${K_p}$: 
\[
\int_{K_p}(\phi ,\phi )_\chi dk<\infty . 
\]

\begin{conjecture}
There are canonically normalized intertwining operators 
\[
R_{P^{\prime }|P}(\chi ):{\cal{K}}_P(\chi )\rightarrow
{\cal{K}}_{P^{\prime}}(\chi ) 
\]
such that

(1) for all $x\in \overline{G({\Bbb{Q}}_p)}$, we have 
\[
R_{P^{\prime }|P}(\chi )I_P(\chi ,x)
=I_P(\chi ,x)R_{P^{\prime }|P}(\chi ), 
\]

(2) $R_{P^{\prime }|P}(\chi )$ is unitary and 
\[
R_{P^{\prime \prime }|P}(\chi )
=R_{P^{\prime \prime }|P^{\prime }}(\chi)R_{P^{\prime }|P}(\chi ), 
\]

(3) if $\psi \in {\cal K}_P(\chi )$ is such that
$\psi (k)=\psi (1)$
for all $k\in \overline{K_p}$ then 
\[
R_{P^{\prime }|P}(\chi )\psi =\psi , 
\]

(4) (for $p=\infty $) if $W\in \Pi (\overline{K_p})$ and 
\[
R_{P^{\prime }|P}(\chi )_W=R_{P^{\prime }|P}(\chi )|_{{\cal{K}}_P(\chi )_W}
\]
then for all $D\in Diff(i{\frak {a}}_P^{*})$ there is a $C>0,\;N>0$ such
that 
\[
||D_\chi R_{P^{\prime }|P}(\chi )_W||\leq C\cdot (1+||\chi ||)^N\cdot
(1+||W||)^N,
\]
for all $\chi \in \Pi _u(\overline{A({\Bbb{Q}}_p)})$.
Here $||\chi||,\;||W||$ are as in \cite{A6}, section 5
(see also section 8.5 above).
\end{conjecture}

If $p=\infty $, $G=SL(2)$ then these follow from \cite{Sa}.
For $p<\infty$ and $n=1$ see \cite{A16}, section 2. If $p<\infty$
and $n=2$, see \cite{GS}. For our immediate concern, it will be
important to choose the normalizing factors for $n=1$
consistent (in a sense to be described later)
with those for $n=2$.

\subsection{The terms in the invariant trace formula}

Recall ${\cal C}$ denotes the set of equivalence classes in $G({\Bbb {Q}})$
(two elements are equivalent if and only if their semi-simple
Jordan components are $G({\Bbb{Q}})$-conjugate). Let
$\sigma \in G({\Bbb{Q}})$ be semi-simple. The
$(G,S)${\bf -equivalence class of }$\sigma $ is
the set 
\[
G({\Bbb{Q}})\cap (\sigma {\cal{U}})^{G({\Bbb{Q}})}
=\{\delta ^{-1}\sigma u\delta
\;|\;\delta \in G({\Bbb{Q}}),\;u\in {\cal{U}}\cap G({\Bbb{Q}})\}, 
\]
where ${\cal{U}}$ is a unipotent conjugacy class in
$\overline{G({\Bbb {A}}_S)}_\sigma$.

\begin{lemma}
Each ${\cal{O}}\in {\cal{C}}$ breaks up into a finite set of
$(G,S)$-classes ${\cal{O}}_{(G,S)}$.
\end{lemma}

More generally, we may define this
finer notion of equivalence for any Levi $M$. Let
$\gamma \in M({\Bbb{Q}})$ have semi-simple Jordan
component $\sigma$. Let $\gamma ^{\prime }\in M({\Bbb{Q}})$.
We say that $\gamma ^{\prime }$
is $(M,S)${\bf -equivalent} to $\gamma $ if there is a
$\delta \in M({\Bbb{Q}})$ such that

(1) $\sigma $ is the semi-simple Jordan component of
$\delta ^{-1}\gamma^{\prime }\delta $,

(2) $\sigma ^{-1}\gamma $ and $\sigma ^{-1}\delta ^{-1}
\gamma ^{\prime}\delta $ are
$\overline{M({\Bbb{A}}_{S})}_\sigma$-conjugate.

\begin{example}
Let $G=SL(2)$, $\gamma,\gamma^{\prime}\in G({\Bbb{Q}})$
and let $S$ be non-empty.

Let $M=G$ and $\gamma =1$. $\gamma$ is $(G,S)$-equivalent to
$\gamma^{\prime}$ if and only if they are equal.

Let $M=A$ and $\gamma =1$. $\gamma$ is $(A,S)$-equivalent to
$\gamma^{\prime}$ if and only if they are equal.

Let $M=G$ and let $\gamma $ be regular semi-simple.
$\gamma$ is $(G,S)$-equivalent to
$\gamma^{\prime}$ if and only if they are $G({\Bbb{Q}})$-conjugate.

Let $M=A$ and $\gamma \in A({\Bbb{Q}})$ be regular.
$\gamma$ is $(A,S)$-equivalent to
$\gamma^{\prime}$ if and only if they are equal.
\end{example}

\begin{lemma}
If $G=SL(2)$ and $S$ is non-empty then
$\gamma^{\prime}$ is $(G,S)$-equivalent to $\gamma\in G({\Bbb{Q}})$ only if 
$\gamma^{\prime}$ is equivalent to $\gamma$.
\end{lemma}

\begin{example}
Let $G=SL(2)$ and let ${\cal{O}}\in {\cal C}-{\cal C}(G)$, so
${\cal{O}}$ is not an elliptic class. If $a$ denotes the Jordan component of
some $\gamma \in {\cal{O}}$ and $a$ is regular then we have
${\cal{O}}\cap A({\Bbb{Q}})=\{a,a^w\}$.
\end{example}

Let $G=SL(2)$. Assume that
${\cal{O}}\in {\cal{C}}-{\cal{C}}(G)$
contains a regular element.
We next define a term $J_{{\cal{O}}}^{\prime }(f)$
such that 
\begin{equation}
I_{{\cal{O}}}(f)=J_{{\cal{O}}}(f)-J_{{\cal{O}}}^{\prime }(f)
\end{equation}
is an invariant distribution on ${\cal{H}}(\overline{G({\Bbb{A}}_S)})$.
To define $J_{\cal{O}}^{\prime }(f)$, let
$\chi \in \Pi _u(\overline{A({\Bbb{A}}_S)})$, let
$M(\chi )=M(w,\chi)$ denote the intertwining operator
\[
M(w,\chi):I_A(\chi )\rightarrow I_A(\chi ^w),
\]
and let $M(\chi ,\eta )$ denote the operator
$M(w,\eta )=M(w,I_A(\chi ),\eta)$ in \cite{A2}, section 4, p. 113,
where $\eta \in {\frak {a}}_{{\Bbb {C}}}^{*},\;P_1=B,\;P_2=B^w$.
Let $m(\chi ,\eta )$ denote the scalar factor such that 
\begin{equation}
R(\chi ,\eta )=m(\chi ,\eta )^{-1}M(\chi ,\eta ):
{\cal K}_B(\chi )\rightarrow {\cal K}_B(\chi ^w)
\end{equation}
satisfies the normalized intertwining operator satisfying the conditions in
section 9.4.4 above. The relation between the logarithmic
derivatives of the two is given by
\begin{equation}
M(\chi ,\eta )^{-1}M'(\chi ,\eta )=
m(\chi ,\eta )^{-1}m'(\chi ,\eta )+
R(\chi ,\eta )^{-1}R'(\chi ,\eta ).
\end{equation}

Let 
\begin{equation}
\varphi _A(f,\chi )
={\rm tr}(R(\chi ,0)^{-1}R(\chi ,0)^{\prime }I_A(\chi ,f)),
\end{equation}
where the derivative (the prime on $R$) is with respect to $\eta $ and then
evaluated at $0$. This makes sense since the range of $R(\chi ,\eta )$ is
independent of $\eta $.
Roughly speaking, the correction term $J_{\cal{O}}^{\prime }(f)$ can be
regarded as an average of the inverse Fourier transform of $\varphi _A(f)$
over ${\cal{O}}\cap A({\Bbb{Q}})$
(by analogy with the case of $SL(2)$ or $GL(2)$
- see the introduction to \cite{A6} and \cite{LL}). This is
essentially the way
it works in the formalism here but more work will be required to set things
up properly.

We now turn to \cite{A6}, section 4 and recall from
chapter 3 above the
definition of the invariant distributions
$I_{{\cal{O}}}^A,\;I_\chi ^A$ on $\overline{A({\Bbb{A}}_S)}$.
These will be used to define $I_{{\cal{O}}}^G,\;I_\chi ^G$ on 
$\overline{G({\Bbb{A}}_S)}$ using the metaplectic analog of
\cite{A6}, Proposition 4.1 and its proof. In fact, we
have 
\begin{equation}
J_{{\cal{O}}}^{\prime }(f)
=|W_G|^{-1}(I_{{\cal{O}}}^A)^{\wedge }(\varphi^G_A(f)),
\end{equation}
where $(I_{{\cal{O}}}^A)^{\wedge }$ denotes the distribution on the
Schwartz space ${\cal{S}}(\Pi _u(A({\Bbb {A}}_S)))$ such that 
\[
(I_{\cal{O}}^A)^{\wedge }(F_A)=I_{\cal{O}}^A(F)\ \ \ \ (=F({\cal{O}})), 
\]
where 
\[
\varphi_A^G:{\cal{H}}(\overline{G({\Bbb {A}}_S)})\rightarrow
I(\Pi _u(\overline{A({\Bbb {A}_S)})}) 
\]
is the weighted character in (9.64) and $F_A$ (a function on
$\Pi_u(\overline{A({\Bbb {A}_S)})}) $) denotes the Fourier transform
of $F\in {\cal{H}}(\overline{A({\Bbb {A}}_S)})$.

\begin{proposition}
(\cite{A6}, Introduction and Proposition 4,1) The distribution
(9.66) on ${\cal{H}}(\overline{G({\Bbb{A}}_S)})$ is invariant,
where $J_{\cal{O}}^{\prime }(f)$ is given
by (9.65).
\end{proposition}

Assume that ${\chi }\in {\cal{X}}-{\cal{X}}(G)$.
Next we define a term $J_{{ \chi }}^{\prime }(f)$ such that 
\begin{equation}
I_{{\chi }}(f)=J_{{\chi }}(f)-J_{{\chi }}^{\prime }(f)
\end{equation}
is an invariant distribution on
${\cal{H}}(\overline{G({\Bbb{A}}_S)})$.
Let 
\begin{equation}
J_{{\chi }}^{\prime }(f)=|W_G|^{-1}(I_{{\cal \chi }}^A)^{\wedge
}(\varphi _A(f)),
\end{equation}
where $(I_{{\cal \chi }}^A)^{\wedge }$ denotes the distribution on the
space $I(\Pi _u(\overline{A({\Bbb {A}}_S)}))$ such that 
\[
(I_{{\cal \chi }}^A)^{\wedge }(F_A)=I_{{\cal \chi }}^A(F). 
\]

\begin{proposition}
(\cite{A6}, Introduction and Proposition 4.1) The distribution
(9.70) on ${\cal{H}}(\overline{G({\Bbb{A}}_S)})$ is invariant,
where $J_{{\cal \chi }}^{\prime }(f)$ is given by (9.67).
\end{proposition}

\subsection{The terms $I_M(\gamma, f)$ and $I_M(\pi,f)$}

The term $I_M(\gamma,f)$ is defined inductively by
\begin{equation}
I_M(\gamma,f)
=J_M(\gamma,f)-
\sum\Sb L\\ M\subset L\\ L\not= G\endSb
\hat{I}_M^L(\gamma,\phi_L(f)),
\end{equation}
where $\phi_L$ is as in 9.4.2.
For a general explicit expression, see \cite{A16}.

There is an ``invariant version'' of Proposition 279, where $J$ is
replaced by $I$: Let $(M({\Bbb{Q}})\cap {\cal{O}})_{M,S}$ denote the set of
$(M,S)$-equivalence classes in $M({\Bbb {Q}})\cap {\cal{O}}$.

\begin{proposition}
(\cite{A10}, Theorem 3.3) 
\[
I_{{\cal{O}}}(f)
=\sum_{M\in {\cal{L}}(G)}|W_0^M|\cdot |W_0^G|^{-1}
\sum_{\gamma \in (M({\Bbb{Q}})\cap {\cal{O}})_{M,S}}
a^M(S,\gamma )I_M(\gamma ,f),
\]
where $S$ is a sufficiently large set of valuations and the
$a^M(S,\gamma)$ are as in section 9.2.8.
\end{proposition}

In section 9.4.4 we introduced normalized intertwining operators
$R_{P^{\prime}|P}(\pi)$, for $P^{\prime},P\in {\cal P}(M)$.
We will use these to make the distributions
$I_M(\gamma)$ and $I_M(\pi)$ more explicit.
For a fixed $P_0\in {\cal P}(M)$,
$\nu\in i{\frak{a}}_M^*$,
$\pi\in \Pi_u(\overline{M({\Bbb{A}}_S)})$,
$\lambda\in {\frak{a}}_{M,{\Bbb{C}}}^*$, let
\begin{equation}
R_P(\mu,\pi_\lambda,P_0)
=R_{P|P_0}(\pi_\lambda)^{-1}R_{P|P_0}(\pi_{\lambda+\mu}) .
\end{equation}
This more than likely (the metaplectic analog of) a
$(G,M)$-family \cite{A11}, section 6. In any case, the quantity
\[
R_M(\pi_\lambda,P_0)
=\lim_{\mu\rightarrow 0}
\sum_{P\in {\cal P}(M)}
R_P(\mu,\pi_\lambda,P_0)\theta_P(\mu)^{-1}
\]
exists. Here $R_G(\pi_\lambda,P_0)=1$ and
\[
R_A(\pi_\lambda,P_0)=
R_{\overline{P_0}|P_0}(\pi_\lambda)^{-1}
{d\over d\lambda}
R_{\overline{P_0}|P_0}(\pi_\lambda)
\]
by the computation in Example 303 (see also \cite{A11} and
\cite{A6}, (6.5)).

\begin{example}
If $G=SL(2)$ and $M=A$ then
\[
I_A(\gamma,f)
=J_A(\gamma,f)-
\hat{I}_A^A(\gamma,\phi_A(f)),
\]
for $\gamma\in A({\Bbb{Q}})$.
The term $\hat{I}_A^A(\gamma,\phi_A(f))$
is the inverse Fourier transform of
$\phi_A(f)(\chi)=tr(I_A(\chi)R_A(\chi,B))(f)$ evaluated at
$\gamma$. 
\end{example}

Let
\begin{equation}
J_M(\pi_\lambda,f)=
{\rm tr}(R_M(\pi_\lambda,P_0)
I_{P_0}(\pi_\lambda,f)),
\end{equation}
where $f\in {\cal{H}}(\overline{G({\Bbb{A}}_S)})$ and where
$I_{P_0}(\pi)$ denotes the usual induced representation.
Define
\begin{equation}
J_M(\pi,X,f)
=\int_{i{\frak{a}}^*_{M,S}}
J_M(\pi_\lambda,f)e^{\lambda(X)}d\lambda,
\end{equation}
where $f$ is in the Hecke algebra as above, $X\in {\frak{a}}_{M,S}$
is as in (9.55), and where

\begin{eqnarray*}
&&{\frak{a}}^*_{M,S}
={\frak{a}}^*_{M}/{\frak{a}}^{\vee}_{M,S}\\
&&{\frak{a}}^{\vee}_{M,S}={\rm Hom}({\frak{a}}_{M,S},2\pi {\Bbb Z}).
\end{eqnarray*}
For
$\pi\in \Pi_u(\overline{M({\Bbb{A}}_S)})$,
inductively define $I_M(\pi,X,f)=I_M^G(\pi,X,f)$ by
\begin{equation}
I_M(\pi,X,f)
=J_M(\pi,X,f)-
\sum\Sb L\\ M\subset L\\ L\not= G\endSb
\hat{I}_M^L(\pi,X,\phi_L(f)),
\end{equation}
where $\phi_L$ is as in \cite{A11}, \S 12 .
In general, define $I_A(\pi,f)=I_A(\pi,0,f)$.

\begin{example}
If $G=SL(2)$ and $M=A$ then
\[
I_A(\pi,X,f)
=J_A(\pi,X,f)-
\hat{I}_A^A(\pi,X,\phi_A(f))
=J_A(\pi,X,f)-(\phi_A(f))(\pi),
\]
for $\pi\in \Pi_u(\overline{A({\Bbb{A}}_S)})$.
\end{example}

The analog of Proposition 320 is

\begin{proposition}
(\cite{A9}, Lemma 3.1) If $\pi\in\Pi_t(\overline{M({\Bbb{A}}_S)})$ then
\[
I_M(\pi,X,f)=
\left\{
\begin{array}{c}
 f_G(\pi,X),\ \ M=G,\\
 0,\ \ M\not= G,
\end{array}
\right.
\]
where
\[
f_G(\pi,X)=
\int_{i{\frak{a}}^*_{G,S}}
{\rm tr}(\pi_\lambda(f))e^{\lambda(X)}d\lambda,
\]
and $X\in {\frak{a}}_{G,S}$.
\end{proposition}

(This vanishing result for $I_M,\ M\not=G$, is false in
general if $\pi$ is not tempered.)

\begin{example} If $G=SL(2)$ and
$\pi\in\Pi_t(\overline{M({\Bbb{A}}_S)})$ then
\[
I_M(\pi,X,f)=
\left\{
\begin{array}{c}
 tr(\pi(f)),\ \ M=G,\\
 0,\ \ M= A.
\end{array}
\right.
\]
Since $\overline{A({\Bbb{A}}_S)}$ is abelian,
$\Pi_t(\overline{M({\Bbb{A}}_S)})=
\Pi_u(\overline{M({\Bbb{A}}_S)})$. Thus
$I_M(\pi,X,f)=0$.
\end{example}

We assume (as usual) that $S$ contains the archimedian place.
If $\pi$ is an automorphic representation of
$\overline{M({\Bbb{A}}_S)^1}$,
the infinitesimal character of the archimedian
component of $I_M(\pi)$, is a $W$-orbit in the set of
${\frak{a}}_{M,{\Bbb{C}}}^*$-orbits in ${\frak{h}}_{\Bbb{C}}^*$.
In this double orbit there is a unique
$W$-orbit in ${\frak{h}}_{\Bbb{C}}^*$ of smallest norm, which we denote by
$\nu_\pi$. (If $M=G$ the $\nu_\pi$ is related to the eigenvalue
of the Laplacian of $\phi_\infty$,
where $\phi =\otimes \phi_p$
belongs to the space of $\pi$.)
If $\chi\in {\cal{X}}$ is the class of the orbit of
$(M,\pi)$, $\pi$ an irreducible cuspidal representation of $M$ then let
$\nu_\chi =\nu_\pi$.

\begin{theorem} (\cite{A10}, Corollary 4.5) 
\[
I_\chi
(f)
=\sum_{M\in L(G)}|W_0^M|\cdot |W_0^G|^{-1}
\int_{\Pi (M,\chi)}a^M(\pi )I_M(\pi ,f)d\pi ,
\]
where 
\[
\Pi (M,\chi )=
\{\pi \in \Pi (\overline{M({\Bbb{A}})^1})\;|\;v_\pi =\nu_\chi \},
\]
$a^M(\pi )$ satisfies (4.5) in \cite{A10} (see also below)
and $I_M$ is defined by (9.72).
\end{theorem}

By the above example, if $G=SL(2)$ and if
$\chi\in {\cal X}-{\cal X}(G)$ then the term $M=A$ in
the above expression for $I_\chi (f)$ is zero. Roughly speaking,
this terms have been put into the ``hyperbolic terms'' in
the geometric expansion.
 
\subsection{The invariant trace formula}

The geometric side of the trace formula is the sum of the following terms:

(a) the sum over classes ${\cal{O}}\subset G({\Bbb{Q}})$
such that ${\cal{O}}\cap Z(G({\Bbb{Q}}))
\not= \emptyset$ of
\begin{equation}
I_{\cal{O}}(f)=\sum_{\gamma\in {\cal{O}}_{(G,S)}}
a^G(S,\gamma)I_G(\gamma,f)
+{1\over 2}\sum_{z\in ({\cal{O}}\cap Z(G({\Bbb{Q}})))_{(A,S)}}
a^A(S,z)I_A(z,f),
\end{equation}
where $I_M$ and $a^M$ are given in (9.72) and section 9.2.8,
resp., and the $(G,S)$-classes were described in Example 315
above (these include the ``nilpotent terms''),

(b) the sum over elliptic regular classes 
${\cal{O}}\subset G({\Bbb{Q}})$
of
\begin{equation}
I_{\cal{O}}(f)=
\sum_{\gamma\in {\cal{O}}_{(G,S)}}
a^G(S,\gamma)J_G(\gamma,f),
\end{equation}
where $\gamma\in {\cal{O}}$ and
$J_G$ (an orbital integral) and $a^G$ (a volume)
are given in section 9.2.6,

(c) the sum over ``hyperbolic'' regular classes
${\cal{O}}\subset G({\Bbb{Q}})$
of
\begin{equation}
I_{\cal{O}}(f)=
\sum_{z\in ({\cal{O}}\cap A({\Bbb{Q}}))_{(A,S)}}
a^A(S,\gamma)I_A(\gamma,f),
\end{equation}
where $I_A$ and $a^A$ (a volume)
are given in above and in section 9.2.6.

The spectral side of the trace formula is the sum of the following terms:

(d) the sum over classes $\chi\in {\cal{X}}(G)$ of
\begin{equation}
I_\chi(f)=
{\rm tr}(\pi(f)),\ \ \ (\pi,G)\in\chi,
\end{equation}

(e) the sum over classes $\chi\in {\cal{X}}-{\cal{X}}(G)$ of
\begin{equation}
I_\chi(f)=J_\chi(f)-J'_\chi(f)
=\int_{\Pi (G,\chi)}a^G(\pi ){\rm tr}\pi (f)d\pi ,
\end{equation}
as in Theorem 325. Here $a^G(I_A(\sigma))=-m'(\sigma)/m(\sigma)$.



